\section{Finite commensurate flat sectors}
\label{sec:flat-sectors}

An exponential that is smaller than every power of the source coordinate
needs its own truncation index. A finite part of this larger scale arises
from the following exact forward expression in a positive source coordinate
$u\to0^+$:
\begin{equation}
 F(u)=y_0+a u^q+\sum_{j=1}^J A_j(u)e^{-k_j c/u^p},
 \quad p,c>0,\quad q\ne0,\quad a\ne0,\quad k_j\in\N_{>0},
 \label{eq:flat-forward}
\end{equation}
where each $A_j$ is a finite real power--log expression. The phase power and
rates are fixed real numbers, with all rates integer multiples of a common
positive rate. Coefficients may depend on fixed real parameters subject to
the stated sign assumptions. Different phase powers and incommensurate
rates require more than one exponential coordinate and fall outside this
single-phase construction.

The zero sector must be exact. An ordinary finite asymptotic approximation
to a nonmonomial core leaves a power-sized error, which dominates every
sector in \eqref{eq:flat-forward}. The exact shifted monomial core
$y_0+a u^q$ avoids this obstruction: its zero-sector inverse is exact. Put
\[
 u_0=\left(\frac{y-y_0}{a}\right)^{1/q}>0,
 \qquad E_0=e^{-c/u_0^p}.
\]
At a finite source endpoint the observable is
$H(u)=x_0+\sigma u$ for power one and $H(u)=\sigma^r u^r$ otherwise.
At an infinite endpoint it is $H(u)=\sigma^r u^{-r}$.
For negative source branches, take an integer observable power.

\subsection{Computing complete sectors}

Introduce a formal symbol $Z$ and the polynomial
$R(u,Z)=\sum_j A_j(u)Z^{k_j}$. Differentiation of the actual exponential
coordinate is represented by the exact operator
\[
 \mathcal D=\partial_u+cp\,u^{-p-1}Z\partial_Z,
 \qquad \frac{d}{du}e^{-c/u^p}=cp\,u^{-p-1}e^{-c/u^p}.
\]
The marker identity of Section~\ref{sec:core-perturbations} now gives
\begin{equation}
 \frac{(-1)^n}{n!}
 \left(\frac{\mathcal D}{a q u^{q-1}}\right)^{n-1}
 \left(\frac{H'(u)R(u,Z)^n}{a q u^{q-1}}\right).
 \label{eq:flat-marker-coefficients}
\end{equation}
Sum this expression over $1\le n\le N$, group equal powers of $Z$, and
evaluate $u=u_0$, $Z=E_0$. Since every $k_j\ge1$, no marker degree larger
than $N$ contributes to an exponential degree at most $N$. Also
$\mathcal D$ preserves exponential degree, so intermediate polynomial
truncation discards no contribution to a retained sector. Each resulting
coefficient is a complete finite power--log function of $u_0$.

For example, the inverse of $F(x)=x+e^{-1/x}$ at $0^+$ has the
first three exponential corrections
\begin{equation}
 y-E+y^{-2}E^2+
 \left(y^{-3}-\frac32 y^{-4}\right)E^3,
 \qquad E=e^{-1/y}.
 \label{eq:flat-primary-inverse}
\end{equation}
For $F(x)=x+x^2e^{-1/x}$, the first two sectors are
$y-y^2E+(y^2+2y^3)E^2$. Adding $2e^{-2/x}$ to the first example changes
the second coefficient to $y^{-2}-2$. Thus the sector index distinguishes
input sectors from powers of the perturbation marker.

\subsection{A bound for the complete omitted tail}

Write every coefficient as a finite sum $u^\alpha P(\log u)$, and choose
\[
\begin{aligned}
 b=\max(0,\max_{\alpha}(q+p-\alpha)),\qquad
 d&=\max(0,\max_P\deg P),\\
 B(u)=u^{-b}\M(u)^d,\qquad \chi(u)=B(u)e^{-c/u^p}.
\end{aligned}
\]
Let $r_*=r$ at a finite source endpoint and $r_*=-r$ at infinity.
For fixed admitted data, there are $C,K>0$ and a source threshold such that
the complete omitted sector tail after degree $N$ satisfies
\begin{equation}
 |\operatorname{Tail}_N|
 \le C u_0^{r_*+p}
 \frac{(K\chi(u_0))^{N+1}}{1-K\chi(u_0)}
 \quad\hbox{when }K\chi(u_0)<1.
 \label{eq:flat-complete-tail}
\end{equation}
Consequently its asymptotic class is
\[
 \OO\!\left(E_0^{N+1}
 u_0^{r_*+p-(N+1)b}\M(u_0)^{(N+1)d}\right).
\]
This bound is allowed to be conservative; it is independent of whether the
first omitted coefficient cancels.

\paragraph{Products.}
Let $s=\sum_{k\le N_a}E_0^kA_k+E_0^{N_a+1}R_a$ and
$t=\sum_{k\le N_b}E_0^kB_k+E_0^{N_b+1}R_b$ with power--log bounds on the
coefficients and tails. The product retains depth $N=\min(N_a,N_b)$, and
every omitted contribution is a graded term $E_0^{k}\OO(u_0^{\beta}\M^{d})$
with $k>N$: the coefficient products $A_iB_j$ with $i+j>N$ at grade $i+j$,
the products $R_aB_j$ and $R_bA_i$ at grades $N_a+1+j$ and $N_b+1+i$, and
$R_aR_b$ at grade $N_a+N_b+2$. For fixed data and $k_2>k_1$,
\[
 \frac{E_0^{k_2}u_0^{\beta_2}\M^{d_2}}{E_0^{k_1}u_0^{\beta_1}\M^{d_1}}
 =\exp\!\Bigl(-\frac{(k_2-k_1)c}{u_0^{p}}\Bigr)\,u_0^{\beta_2-\beta_1}\M^{d_2-d_1}
 \longrightarrow0,
\]
so the complete omitted tail is bounded by the candidates of least grade
alone, combined by power--log dominance, and may then be weakened to grade
$N+1$ for the representation above. Combining the algebraic pairs of all
candidates before comparing grades is also valid but can be much weaker:
for the depth-$N$ inverse of $y=x+e^{-1/x}$ squared, the dominant candidate
is $R_a A_0$ at grade $N+1$ with power $1-2N$, whereas the tail product
$R_aR_b$ at grade $2N+2$ has power $-4N$ and would otherwise control the
bound. Exact cancellations among the coefficient products of the first
omitted grade can be found by multiplying them out; deeper omitted grades
need only their envelopes.

To justify the estimate, set $u=u_0(1+u_0^p V)$ and divide the forward
equation by $u_0^{q+p}$. The normalized monomial difference has the analytic
extension
\[
 a\frac{(1+\epsilon V)^q-1}{\epsilon},
 \qquad \epsilon=u_0^p,
\]
whose derivative at $V=0$, $\epsilon=0$ is $a q\ne0$. Each exponential
factor, after extracting $E_0^{k_j}$, becomes
\[
 \exp\!\left(k_jc\,
 \frac{1-(1+\epsilon V)^{-p}}{\epsilon}\right),
\]
which is analytic and uniformly bounded for small $\epsilon,V$. The
power--log factors have the same local analytic property after their
powers of $u_0$ and $\log u_0$ are extracted. Each normalized coefficient
is bounded by a constant times $B(u_0)$. Since $B\ge1$ for small $u_0$,
replacing $E_0$ by a complex variable $Z$ gives a uniform implicit solution
on $|Z|<1/(K B(u_0))$. One may regard the bounded normalized coefficients
as additional parameters in a compact polydisc; the derivative condition
at $Z=0$ is uniform. Cauchy estimates therefore bound the coefficient of
$Z^k$ by $C(KB)^k$. The nonconstant observable contributes the factor
$u_0^{r_*+p}$. Summing the geometric tail proves
\eqref{eq:flat-complete-tail}, and $\chi(u_0)\to0$ establishes applicability
at $Z=E_0$.

The constants and threshold in \eqref{eq:flat-complete-tail} are
existential. They are not explicit numerical error bounds until suitable
values have been established. Sector depth is inclusive: a depth $N$
retains $k\le N$, whereas an ordinary power cutoff retains powers strictly
below its boundary. The first omitted coefficient and a bound for the
complete omitted tail are distinct mathematical objects.

The primary example also admits a direct formal check. Substitute
$g=y+\sum_{k=1}^N C_k(y)Z^k$ into
\[
 g+Z\exp(1/y-1/g)-y.
\]
Vanishing through degree $N$ in $Z$ identifies the same sector coefficients
by the uniqueness of the formal implicit solution. This identity verifies
the finite coefficients; the uniform analytic argument above supplies the
bound for the complete infinite tail.
